% §1 Introduction — to be REWRITTEN, not transcribed (draft §1 is too dense: each of its
% paragraphs 3–7 compresses a whole section).
%
% Shape (2005 §1, four moves): conceptual hook → prior work → deficiency, one line each →
% roadmap. Decisions of record:
%   - New-results-led, positive-first (ssf decision, 2026-07-17): open on the hemigroup
%     move and the self-decomposable class with its finite-moment members; the 2005 heavy
%     tails are stated as fact, the point of departure, not the frame.
%   - One-paragraph Lindeberg positioning: the time-causal limit kernel is the
%     first-discovered member of the class characterized here (containment, not
%     competition — the content of draft Remark 15.1, one paragraph of it up front).
%   - Embodiment recurs as theorem content TWICE (signaling form §9, with the N-jet §12
%     as its consequence); Sonine conservation is a further consequence, in Appendix A —
%     demoted from the draft's three-instance claim (decision 2026-08-12).
%   - §1.1 "What is machine-checked" = draft §1.5, essentially as-is.
%
% Source: draft/hemigroup-causal-scale-space-kernels.md lines 5–55.

\section{Introduction}\label{sec:introduction}

% [to be written — see task list; drafted after §2 so the two do not overlap]

\subsection{What is machine-checked}\label{sec:machine-checked}

% [draft §1 "What is machine-checked", lines 31–55, essentially verbatim]
